\chapter{Gestión económica}
\label{chap:gestion-economica}

\begingroup
\normalsize
\setstretch{1.05}
\setlength{\parskip}{0.3\baselineskip}
\newcolumntype{Y}{>{\raggedright\arraybackslash}X}
\newcolumntype{R}[1]{>{\raggedleft\arraybackslash}p{#1}}
\newcolumntype{L}[1]{>{\raggedright\arraybackslash}p{#1}}

El presupuesto parte de la planificación temporal: 540 horas repartidas en las
tareas del diagrama de Gantt y los recursos R1--R6. Sigue la estructura
propuesta en GEP: costes de personal por actividad (CPA), costes genéricos
(CG), contingencia e imprevistos. Todos los importes se expresan en euros.

\section{Costes de personal por actividad}
\label{sec:cpa}

El estudiante asume los cuatro roles de la planificación. El coste horario de
cada rol corresponde a un perfil de recién graduado, ya que el TFG lo realiza una
persona que se incorpora al mercado laboral. Por ello, se toma el mínimo del
rango salarial habitual que publica Indeed España (septiembre de 2026) para el
perfil más próximo, en lugar de la media. El salario bruto anual se divide entre 1\,760 horas
anuales (220 días laborables de 8 horas) y se multiplica por 1,3 para incluir
la cotización a la Seguridad Social.

\begin{table}[H]
  \centering
  \small
  \begin{tabularx}{\textwidth}{L{2.9cm}YR{1.6cm}R{1.3cm}R{1.5cm}}
    \toprule
    \textbf{Rol} & \textbf{Perfil de referencia} & \textbf{Bruto anual (mín.)} &
    \textbf{Bruto/h} & \textbf{Coste/h (×1,3)} \\
    \midrule
    R1-A Jefe de proyecto & Jefe de proyecto \citep{SalarioJefeProyecto} & 25\,284\,€ & 14,37\,€ & 18,68\,€ \\
    R1-B Investigador & Científico de datos \citep{SalarioDataScientist} & 27\,671\,€ & 15,72\,€ & 20,44\,€ \\
    R1-C Desarrollador & Desarrollador de software \citep{SalarioDesarrolladorSoftware} & 22\,262\,€ & 12,65\,€ & 16,44\,€ \\
    R1-D Analista de resultados & Analista de datos \citep{SalarioAnalistasDatos} & 21\,903\,€ & 12,44\,€ & 16,18\,€ \\
    \bottomrule
  \end{tabularx}
  \caption[Coste horario por rol]{Coste horario por rol, incluida la Seguridad
  Social, para un perfil recién graduado. Elaboración propia a partir de Indeed España.}
  \label{tab:coste-roles}
\end{table}

La tabla~\ref{tab:cpa} asigna las horas de cada tarea a los roles indicados en
la planificación. Cuando una tarea combina roles, las horas se reparten según
la parte del trabajo que corresponde a cada uno. La redacción de la memoria
(GEP05) se imputa al investigador, porque consiste en documentar el método,
los experimentos y sus resultados. Las horas del director y
del codirector no se presupuestan, porque no forman parte de las 540 horas del
estudiante.

\begingroup
  \footnotesize
  \setlength{\tabcolsep}{3pt}
  \renewcommand{\arraystretch}{0.95}
  \setlength{\LTleft}{0pt}
  \setlength{\LTright}{0pt}
  \begin{longtable}{L{1.1cm}L{5.5cm}R{0.7cm}R{0.7cm}R{0.7cm}R{0.7cm}R{0.7cm}R{2.0cm}}
    \toprule
    \textbf{Cód.} & \textbf{Tarea} & \textbf{A} & \textbf{B} & \textbf{C} &
    \textbf{D} & \textbf{h} & \textbf{Coste} \\
    \midrule
    \endfirsthead
    \multicolumn{8}{c}{\tablename\ \thetable\ (continuación)} \\
    \toprule
    \textbf{Cód.} & \textbf{Tarea} & \textbf{A} & \textbf{B} & \textbf{C} &
    \textbf{D} & \textbf{h} & \textbf{Coste} \\
    \midrule
    \endhead
    \midrule
    \multicolumn{8}{r}{\small Continúa en la página siguiente} \\
    \endfoot
    \bottomrule
    \caption[Costes de personal por actividad]{Costes de personal por actividad
    del Gantt. A--D son las horas de los roles R1-A--R1-D. Elaboración
    propia.}\label{tab:cpa}\\
    \endlastfoot
    EP01 & Revisar el estado del arte & & 28 & & & 28 & 572,32\,€ \\
    EP02 & Analizar los datos y requisitos & & 20 & & & 20 & 408,80\,€ \\
    EP03 & Definir el protocolo experimental & 2 & 8 & & 4 & 14 & 265,60\,€ \\
    EP04 & Estudiar alternativas de despliegue & & 6 & 4 & & 10 & 188,40\,€ \\
    \midrule
    & \textbf{Subtotal Estudio Previo} & 2 & 62 & 4 & 4 & 72 & \textbf{1\,435,12\,€} \\
    \midrule
    GEP01 & Preparar la entrega 1 & 6 & & & & 6 & 112,08\,€ \\
    GEP02 & Preparar la entrega 2 & 6 & & & & 6 & 112,08\,€ \\
    GEP03 & Preparar la entrega 3 & 5 & & & & 5 & 93,40\,€ \\
    GEP04 & Realizar el seguimiento y control & 7 & & & & 7 & 130,76\,€ \\
    GEP05 & Redactar la memoria & & 60 & & & 60 & 1\,226,40\,€ \\
    GEP06 & Preparar y realizar la defensa & 4 & & & 2 & 6 & 107,08\,€ \\
    \midrule
    & \textbf{Subtotal Gestión del Proyecto} & 28 & 60 & & 2 & 90 & \textbf{1\,781,80\,€} \\
    \midrule
    DT01 & Preparar el entorno reproducible & & & 12 & & 12 & 197,28\,€ \\
    DT02 & Adquirir y validar los datos & & 4 & 14 & & 18 & 311,92\,€ \\
    DT03 & Construir el flujo de datos & & & 28 & & 28 & 460,32\,€ \\
    DT04 & Implementar particiones y unificar datos & & 4 & 10 & & 14 & 246,16\,€ \\
    DT05 & Implementar y ejecutar la línea base & & & 14 & & 14 & 230,16\,€ \\
    DT06 & Implementar la CNN & & & 24 & & 24 & 394,56\,€ \\
    DT07 & Entrenar y ajustar la CNN & & & 34 & & 34 & 558,96\,€ \\
    DT08 & Adaptar MobileNet & & & 24 & & 24 & 394,56\,€ \\
    DT09 & Entrenar y ajustar MobileNet & & & 34 & & 34 & 558,96\,€ \\
    DT10 & Adaptar un transformador visual & & & 26 & & 26 & 427,44\,€ \\
    DT11 & Entrenar y ajustar el transformador & & & 36 & & 36 & 591,84\,€ \\
    DT12 & Convertir y cuantizar los modelos & & & 20 & & 20 & 328,80\,€ \\
    DT13 & Validar la ejecución \emph{edge} & & & 12 & 4 & 16 & 262,00\,€ \\
    \midrule
    & \textbf{Subtotal Desarrollo Técnico} & & 8 & 288 & 4 & 300 & \textbf{4\,962,96\,€} \\
    \midrule
    AR01 & Aplicar la evaluación común & & & 14 & 14 & 28 & 456,68\,€ \\
    AR02 & Analizar errores y robustez & & & & 18 & 18 & 291,24\,€ \\
    AR03 & Sintetizar la comparación & & & & 20 & 20 & 323,60\,€ \\
    AR04 & Analizar sostenibilidad y ética & 4 & & & 8 & 12 & 204,16\,€ \\
    \midrule
    & \textbf{Subtotal Análisis de Resultados} & 4 & & 14 & 60 & 78 & \textbf{1\,275,68\,€} \\
    \midrule
    & \textbf{Total CPA} & 34 & 130 & 306 & 70 & 540 & \textbf{9\,455,56\,€} \\
  \end{longtable}
\endgroup

\section{Costes genéricos}
\label{sec:cg}

Los costes genéricos no se asignan a una tarea concreta.

\begin{description}[leftmargin=0.45cm,labelindent=0cm,style=nextline,
  itemsep=0.2\baselineskip]
  \item[Amortización del portátil (R3).] MSI Modern B7M de 750\,€,
  amortizado en 4 años de 1\,760 horas de uso:
  $750 / (4 \times 1\,760) = 0{,}1065$\,€/h. Durante las 540 horas del
  proyecto, $0{,}1065 \times 540 = 57{,}53$\,€.
  \item[Software (R5, R6).] Python, PyTorch, timm, LiteRT, Git, LaTeX y
  VS Code son gratuitos, y Notion se usa con su plan gratuito: 0\,€.
  \item[Cómputo externo (R4).] Suscripción a Google Colab Pro durante cuatro
  meses, a $\approx$10\,€/mes: 40,00\,€. El importe exacto se confirmará al cerrar el proyecto.
  \item[Espacio de trabajo.] Se imputa el coste de un puesto en un
  \emph{coworking} de Barcelona. El pase de día tiene una mediana de 20\,€
  \citep{OnecoworkingcomCitiesBarcelona}; con una jornada de 8 horas, 2,50\,€/h.
  Para las 540 horas del proyecto, $2{,}50 \times 540 = 1\,350{,}00$\,€. El
  pase incluye conexión a Internet y electricidad, por lo que estos conceptos
  no se imputan por separado.
  \item[Dispositivo \emph{edge} (R5).] Está por concretar: puede ser un teléfono
  móvil o una placa como una Raspberry Pi. Se estima un coste de 350\,€ y se imputa
  completo, ya que se adquiere solo para la validación \emph{edge} del proyecto.
\end{description}

\begin{table}[H]
  \centering
  \small
  \begin{tabular}{lr}
    \toprule
    \textbf{Concepto} & \textbf{Coste} \\
    \midrule
    Amortización del portátil & 57,53\,€ \\
    Software & 0,00\,€ \\
    Google Colab Pro & 40,00\,€ \\
    Espacio de trabajo (\emph{coworking}) & 1\,350,00\,€ \\
    Dispositivo \emph{edge} (estimación) & 350,00\,€ \\
    \midrule
    \textbf{Total CG} & \textbf{1\,797,53\,€} \\
    \bottomrule
  \end{tabular}
  \caption[Costes genéricos]{Costes genéricos. Elaboración propia.}
  \label{tab:cg}
\end{table}

\section{Contingencia}

Se aplica una contingencia del 15\,\%, dentro del rango habitual del 10--20\,\%
para proyectos informáticos. Las horas de la planificación ya incluyen un
margen por tarea y el presupuesto está detallado por actividad, por lo que no
se justifica el extremo superior:
$(9\,455{,}56 + 1\,797{,}53) \times 0{,}15 = 1\,687{,}96$\,€.

\section{Imprevistos}

Los imprevistos proceden de los riesgos identificados en la planificación. Cada
uno se imputa según su probabilidad estimada.

\begin{table}[H]
  \centering
  \small
  \begin{tabularx}{\textwidth}{YR{1.5cm}R{1.2cm}R{1.7cm}}
    \toprule
    \textbf{Imprevisto y plan alternativo} & \textbf{Coste} &
    \textbf{Prob.} & \textbf{Imputado} \\
    \midrule
    Cómputo insuficiente en Colab Pro: migrar los entrenamientos a RunPod.
    $\approx$150\,h de GPU (18 configuraciones a $\approx$8\,h) a 0,74\,\$/h
    (RTX 4090, \emph{Secure Cloud}) \citep{GPUCloudPricing}, con
    1\,\$\,$\approx$\,0,90\,€. & 99,90\,€ & 30\,\% & 29,97\,€ \\
    Avería del portátil: sustituirlo por un equipo equivalente. & 750,00\,€ &
    5\,\% & 37,50\,€ \\
    \midrule
    \textbf{Total imprevistos} & & & \textbf{67,47\,€} \\
    \bottomrule
  \end{tabularx}
  \caption[Imprevistos]{Imprevistos imputados según su probabilidad.
  Elaboración propia.}
  \label{tab:imprevistos}
\end{table}

Los demás riesgos de la planificación (dificultades técnicas, error elevado e
incompatibilidad \emph{edge}) no añaden coste: su exceso de horas se compensa
reduciendo el alcance opcional, sin superar las 540 horas.

\section{Presupuesto total}

\begin{table}[H]
  \centering
  \small
  \begin{tabular}{lr}
    \toprule
    \textbf{Concepto} & \textbf{Coste} \\
    \midrule
    Costes de personal por actividad (CPA) & 9\,455,56\,€ \\
    Costes genéricos (CG) & 1\,797,53\,€ \\
    Contingencia (15\,\%) & 1\,687,96\,€ \\
    Imprevistos & 67,47\,€ \\
    \midrule
    \textbf{Total} & \textbf{13\,008,52\,€} \\
    \bottomrule
  \end{tabular}
  \caption[Presupuesto total]{Presupuesto total del proyecto. Elaboración
  propia.}
  \label{tab:presupuesto-total}
\end{table}

El personal representa el 73\,\% del total, un porcentaje habitual en
proyectos centrados en software.

\section{Control de gestión}
\label{sec:control-gestion}

El control reutiliza el registro semanal de la planificación: horas reales y
horas restantes por tarea. Además, se registrarán las horas de cómputo
consumidas por los entrenamientos de los modelos y se guardarán las facturas de
Colab o de cualquier otro proveedor de cómputo en la nube. Con estos datos se
calculan los siguientes indicadores al cerrar cada tarea y en cada hito:

\begin{itemize}
  \item \textbf{Desviación por horas}: (horas estimadas $-$ horas reales)
  $\times$ coste por hora estimado. Indica si una tarea ha requerido más o
  menos horas de las previstas.
  \item \textbf{Desviación por precio}: (coste por hora estimado $-$ coste por
  hora real) $\times$ horas reales. En el personal es siempre cero, porque el
  coste por hora es fijo; se aplica al cómputo, cuyo precio por hora de GPU
  puede cambiar.
  \item \textbf{Desviación total}: coste presupuestado $-$ coste real, tanto de
  los costes de personal como de los costes genéricos.
  \item \textbf{Uso de las reservas}: porcentaje gastado de la contingencia y
  de la partida de imprevistos.
\end{itemize}

Un valor negativo significa que se ha gastado más de lo previsto. Si una tarea
se desvía más del 15\,\% o el sobrecoste acumulado supera la mitad de la
contingencia, se revisan la planificación y el alcance con la dirección. Si el
sobrecoste se debe a uno de los imprevistos presupuestados, se cubre con su partida; en otro
caso, con la contingencia.

\endgroup
